\begin{abstract}
Speaker-attributed transcripts are only useful when they preserve who said what, but multi-person meetings make speaker labels fragile under overlap, short turns, and noisy diarization. Prior work shows that language models can help speaker attribution when they are fine-tuned or paired with strong acoustic evidence. We ask a narrower practical question: can a lightweight API-based LLM correction layer improve noisy speaker labels while preserving transcript words? We build a controlled pilot from local \amisdm and \simsamu audio. \asrmodel transcribes 90 annotated speech segments, seeded corruption produces noisy speaker labels, and \relabelmodel relabels the same transcript under no-profile and profile-conditioned prompts. Across 30 window/noise cases, no-profile LLM relabeling reduces mean \wder from 0.390 to 0.375, a 3.8\% relative reduction, but the effect is not statistically significant after Holm correction (Wilcoxon one-sided $p=0.137$, Holm $p=0.410$). Profile conditioning is weaker, reducing mean \wder to 0.381. The LLM responses are structurally reliable, with 60/60 valid JSON outputs, but they are conservative: no-profile relabeling leaves labels unchanged in 26/30 cases. The improvement appears only in the two-speaker \simsamu windows; four-speaker \amisdm windows show no LLM gain. These results support LLMs as safe transcript-preserving relabelers, but not as a standalone fix for many-speaker meeting identification.
\end{abstract}
